\documentclass[a4paper,12pt]{ctexart}
\usepackage{xcolor,graphicx,geometry,array,hyperref,booktabs}
\hypersetup{colorlinks=true,linkcolor=blue!70!black}
\geometry{top=1.8cm,bottom=1.8cm,left=1.8cm,right=1.8cm}
\linespread{1.35}
\newenvironment{thinkbox}{\vspace{0.3cm}\begin{center}\begin{minipage}{0.88\textwidth}\colorbox{blue!5}{\begin{minipage}{0.94\textwidth}\vspace{0.15cm}\textbf{【想一想】}}\vspace{0.15cm}\end{minipage}\end{minipage}\end{center}\vspace{0.3cm}}{}
\newenvironment{funfact}{\vspace{0.2cm}\begin{center}\begin{minipage}{0.88\textwidth}\fbox{\begin{minipage}{0.92\textwidth}\vspace{0.1cm}\textbf{【有趣的小知识】}}\vspace{0.1cm}\end{minipage}\end{minipage}\end{center}\vspace{0.2cm}}{}
\newenvironment{figplaceholder}[1]{\vspace{0.2cm}\begin{center}\fbox{\begin{minipage}{0.82\textwidth}\centering\textbf{【插图位置】}\\[0.2cm]#1\end{minipage}}\end{center}\vspace{0.2cm}}{}

\title{\Huge\textbf{法学小故事}}
\date{}
\begin{document}\maketitle\thispagestyle{empty}

\section{为什么要有红绿灯？——规则如何保护所有人}

想象一下：一个没有红绿灯的十字路口。四个方向的车都想先过去——每个人都想"我先走"。结果呢？全都堵在路口中间——谁也走不了。这还不是最坏的情况——如果没有人愿意停，结果就是撞车。

红绿灯就是一种"规则"——它告诉每个人，什么时候该走、什么时候该停。它让所有人的行为变得可预测。法律的作用和红绿灯完全一样——它让社会中每个人的行为有了一个共同的"预期"。

\begin{figplaceholder}
插图：没有红绿灯的十字路口（混乱） vs 有红绿灯（有序）——两个场景对比
\end{figplaceholder}

\textbf{法律的三个层次：}
\begin{itemize}
  \item \textbf{民法和合同法：}你和别人签了一份合同（比如租房合同）——如果对方不按合同办，你去法院起诉——法院根据合同法来判决。法律提供了一个"解决纠纷"的框架
  \item \textbf{行政法和交通法：}你闯红灯了——交通警察开罚单。这是政府对个人行为的"管理"。法律给了政府管理公共事务的权力
  \item \textbf{刑法：}有人抢劫、杀人——公安机关调查、检察院起诉、法院判刑。这是国家对最严重违法行为的"惩罚"——因为只有惩罚才能阻止人们做这样的事
\end{itemize}

\begin{thinkbox}
如果一条法律没有人"执行"（比如一个地方从来不查闯红灯）——那么这条法律还有效吗？在书本上有效，但在实际生活中"失效"了。法律的生命在于执行。
\end{thinkbox}

\begin{funfact}
"法"字的起源：古代中国的"法"（灋）字由三个部分构成——"氵"（水，象征公平，水平面一样平）、"廌"（一种独角神兽——獬豸，据说它能分辨善恶，用角顶坏人）、"去"（去除不公）。这个字本身就完美地概括了法的精神——公正和惩罚。
\end{funfact}

\newpage

\section{包青天的故事——古代法官怎么判案}

北宋时期有一个著名的法官——包拯（999-1062年）。他在民间被称为"包青天"——因为他断案公平，不畏权贵。他做过开封府的知府——相当于今天的北京市市长兼最高法院法官。

\textbf{包拯怎么判案？}

\begin{itemize}
  \item \textbf{证据：}包拯非常重视物证和证人证词——他在没有DNA、没有监控的年代，靠走访调查、收集证物做出判决
  \item \textbf{推理：}有一个著名案例——牛舌案。有人割掉了农民家牛的舌头，农民来告状。包拯说："你回去把牛杀了卖肉。"过了几天，有人来举报——说"那个农民私自宰杀耕牛"。包拯下令抓住了举报者——审问后承认了割牛舌。为什么？因为包拯断定——割牛舌的人和举报杀牛的人是同一个人——否则谁会那么"好心"来举报？这就是经典的"不打自招"。
  \item \textbf{不畏权贵：}包拯最出名的品质是——不管对方是皇亲国戚还是高官——犯了法一样判。《宋史》记载他"立朝刚毅，贵戚宦官为之敛手"——他站在朝堂上的时候，平时跋扈的宦官都不敢乱动
\end{complex}

\textbf{古代法庭 vs 现代法庭：}
古代的法庭相对简单——法官一个人说了算（没有律师/没有陪审团/没有上诉——包拯判了就判了）。现代的法庭分工更细：法官（主持审判、适用法律）、检察官（代表国家起诉）、律师（为被告辩护）、陪审团（决定事实问题）。

\begin{figplaceholder}
插图：对比——左半（包拯在公堂上断案）vs 右半（现代法庭——法官/检察官/律师/被告各有位置）
\end{figplaceholder}

\begin{thinkbox}
如果包拯穿越到今天，他还能不能当法官？——他断案公正、善于推理，这些品质永远不会过时。但他可能需要学习现代法律知识——今天的法律体系比宋朝复杂得多（光刑法就有500多条）。
\end{thinkbox}

\begin{funfact}
包拯的"名言"流传至今——但大部分不是他说的。比如"王子犯法，与庶民同罪"——这句话的意思是"法律面前人人平等"——虽然现实中不一定总能做到，但它确立了一个法制的最高原则：没有人在法律之上。
\end{funfact}

\newpage

\section{十二怒汉——陪审团是怎么工作的？}

1957年的经典电影《十二怒汉》（12 Angry Men）讲了一个故事：一个18岁的男孩被指控杀死了自己的父亲——目击证人作证（"我看见他跑下了楼"）、凶器（一把折叠刀）在男孩家中找到、邻居作证（"我听到了争吵声"）。看起来铁证如山。

但12名陪审员——12个普通人——必须"一致同意"才能定罪，其中只有一个人认为："我还不确定。"

这个人向其他11人逐一提出质疑：
\begin{itemize}
  \item "目击证人说他在15秒内透过一列经过的火车看到了男孩跑下楼——但火车经过时的噪音和震动真的能看清吗？"
  \item "凶器是一把非常特别的折叠刀——男孩确实买了同样的刀，但街上到处都有卖"
  \item "邻居的证词——她隔着一条街、在凌晨1点、在半睡半醒中——能确实地"听到"男孩说'我要杀了你'？"
\end{itemize}

经过激烈辩论，陪审团从11:1（有罪）变成了0:12（无罪）。最后发现——那个男孩确实是清白的。

\textbf{这个电影说明了一个重要的法律原则：}
\begin{quote}
\textbf{"合理怀疑"（Beyond a Reasonable Doubt）}——在刑事审判中，如果存在合理的怀疑，就不能定罪。因为"冤枉一个无辜者"比"放走一个有罪者"的代价更大。用法律术语说：定罪的标准是"排除合理怀疑"（不是"绝对确定"——那不可能）。

这个标准比民事诉讼的标准（"优势证据"——哪一方的证据更可信？）要严格得多。为什么？——因为刑事责任的后果是自由（甚至生命），民事责任只是赔偿金钱。
\end{quote}

\begin{figplaceholder}
插图：陪审团在法庭中的位置——12个人坐在侧面的陪审席，法官在中央，检察官和律师分别在两侧
\end{figplaceholder}

\begin{thinkbox}
你觉得"排除合理怀疑"这个标准对吗？有没有可能因为标准太高"放走"了一个真正的罪犯？——会的。这就是司法体系必须选择的风险：是冒风险冤枉一个无辜者，还是冒风险放走一个有罪者。所有文明国家都选择了前者。
\end{thinkbox}

\begin{funfact}
中国的司法体系没有陪审团制，但有\textbf{人民陪审员制度}——在法院审判案件中，由法院随机抽取普通公民担任人民陪审员。和陪审团不同的是，人民陪审员不是单独决定事实问题——他们和法官一起坐在审判席上，共同参与案件的全部审理过程。目前全国约有20多万名人民陪审员。
\end{funfact}

\end{document}
